% TOP_PROBLEM: {"id": 30006676, "title": "Generalized Riemann Hypothesis for Dirichlet L-functions", "category_id": 1, "category": {"id": 1, "name": "number_theory", "display_name": "Number Theory", "description": "Properties of integers, prime numbers, Diophantine equations.", "slug": "number-theory", "order_index": 1, "created_at": "2026-07-31T15:26:25.670Z"}, "status": "open"}
% FIRST_POSTED: null
% !TeX program = lualatex
% Compile twice: lualatex open_problems_500.tex
% All references and prose are embedded; no BibTeX database or external inputs.
\documentclass[11pt,a4paper]{article}
\usepackage[margin=24mm,headheight=14pt]{geometry}
\usepackage{fontspec}
\setmainfont{Latin Modern Roman}
\setsansfont{Latin Modern Sans}
\setmonofont{Latin Modern Mono}
\usepackage{amsmath,amssymb,mathtools}
\usepackage{unicode-math}
\setmathfont{Latin Modern Math}
\usepackage{microtype}
\usepackage{enumitem,xurl,needspace}
\usepackage[dvipsnames]{xcolor}
\usepackage{hyperref}
\hypersetup{unicode=true,colorlinks=true,linkcolor=MidnightBlue,urlcolor=MidnightBlue,
 pdftitle={Mathematical Research Notebook},pdfauthor={Source-annotated compilation},bookmarksnumbered=true}
\usepackage{bookmark}
\usepackage{tocloft}
\setlength{\cftsecnumwidth}{3em}
\cftsetpnumwidth{2.5em}
\cftsetrmarg{4em}
\usepackage{titlesec}
\titleformat{\section}{\Large\bfseries\raggedright}{\thesection}{0.7em}{}
\usepackage{fancyhdr}
\pagestyle{fancy}\fancyhf{}
\fancyhead[L]{\small\sffamily Open Problems Math | Research notebook}
\fancyhead[R]{\small 27 September 2026}
\fancyfoot[C]{\thepage}
\setlength{\parindent}{0pt}
\setlength{\parskip}{4pt plus 1pt}
\setlength{\emergencystretch}{3em}
\setcounter{tocdepth}{1}
\setcounter{secnumdepth}{1}
\setlist[itemize]{leftmargin=3em,itemsep=3pt,topsep=4pt,parsep=0pt}
\allowdisplaybreaks
\newcommand{\N}{\mathbb{N}}
\newcommand{\Z}{\mathbb{Z}}
\newcommand{\Q}{\mathbb{Q}}
\newcommand{\R}{\mathbb{R}}
\newcommand{\C}{\mathbb{C}}
\newcommand{\F}{\mathbb{F}}
\newcommand{\A}{\mathbb{A}}
\newcommand{\PP}{\mathbb P}
\newcommand{\E}{\mathbb{E}}
\newcommand{\eps}{\varepsilon}
\newcommand{\dd}{\,\mathrm d}
\newcommand{\sref}[2]{\hyperlink{source:#1:#2}{[S#2]}}
\newcommand{\eref}[2]{\hyperlink{extra:#1:#2}{[E#2]}}
% Turn the next switch off for a shorter reading copy. The quotations preserve
% every exactTarget field, including qualifications omitted from a short summary.
\newif\ifshowcatalogue
\showcataloguetrue
\newcommand{\cataloguescope}[1]{\ifshowcatalogue
 \paragraph{Exact scope in the supplied catalogue.}
 \begin{quote}\small #1\end{quote}\fi}
\usepackage{etoolbox}
\AtBeginEnvironment{itemize}{\small}
% Standalone entry; original section content is delimited for verification.
\begin{document}
\setcounter{section}{0}
%% BEGIN ORIGINAL SECTION CONTAINER
% ================================================================
% problemId: 30006676
% Canonical problem ID: 30006676
% exactTarget SHA-256: ba13df3a56ab0f77bb9fc819089774dfd9448270f40673fb9fc9203872d06556
\clearpage
\section*{\texorpdfstring{Generalized Riemann Hypothesis for Dirichlet L-functions}{Generalized Riemann Hypothesis for Dirichlet L-functions}}\label{30006676}
\subsection{Definitions and mathematical statement}
A Dirichlet character modulo $q$ is a homomorphism $(\Z/q\Z)^\times\to\C^\times$, extended periodically to integers and by zero on nonunits. Initially, for $\Re s>1$,
\[
 L(s,\chi)=\sum_{n\ge1}\frac{\chi(n)}{n^s}=\prod_p(1-\chi(p)p^{-s})^{-1}.
\]
Take the meromorphic continuation. For every $q\ge1$ and every character, ask whether
\[
 L(\rho,\chi)=0,\quad0<\Re\rho<1\quad\Longrightarrow\quad\Re\rho=\tfrac12.
\]
Principal and imprimitive characters remain in scope. The restriction to the open strip excludes poles and zeros on its boundary caused by removed Euler factors.
\par\smallskip\textit{Formulation and concept references:} \sref{30006676}{1}.
\cataloguescope{For every integer q>=1 and every Dirichlet character chi modulo q, including principal and imprimitive characters, every zero s of the meromorphic continuation of L(s,chi)=sum\_\{n>=1\}chi(n)n\textasciicircum{}(-s) in 0<Re(s)<1 satisfies Re(s)=1/2. The continuation is from Re(s)>1. Only open-strip zeros are quantified: principal-character poles and boundary zeros of imprimitive Euler factors are not counterexamples. No simplicity or central nonvanishing requirement is added.}
\subsection{Short English statement}
Do all Dirichlet $L$-functions place their zeros inside the critical strip on its central vertical line?
%% BEGIN RESEARCH INSERTION
\subsection{Research attempt}
\paragraph{Current frontier.}
\textit{Research check: 27 September 2026.}
The primitive functional equation, induced-character Euler factors and
open-strip zero conventions are recorded in DLMF Section 25.15.
\sref{30006676}{1}. The August 2026 paper on large quadratic character sums
assumes GRH for the Omega results described in its abstract; it cannot
be used as an unconditional proof of the assumption. \sref{30006676}{2}.
We pursue cancellation in the coefficients of $1/L$, rather than
positivity of a product of $L$-functions. The integral criterion below
is a sufficient condition whose implication is proved here; no new
unconditional cancellation estimate is claimed.

\paragraph{Attack route.}
A spectral model would need a self-adjoint operator with the actual
completed $L$-function as its determinant. A direct character-sum argument
controls $L$ but not necessarily its reciprocal. The chosen route is to
seek a mean-square estimate for Mobius-twisted sums; it would give a
holomorphic reciprocal across the right half of the critical strip.
The missing lemma is quantitative cancellation in those sums, with
careful attention to the principal and imprimitive cases.

\paragraph{Attempt: reduce to primitive characters without losing cases.}
Let $\chi\bmod q$ be induced by a primitive character $\chi^*$ of
conductor $q^*\mid q$. Euler products in $\Re s>1$ give
\[
 L(s,\chi)=L(s,\chi^*)
       \prod_{\substack{p\mid q\\p\nmid q^*}}
                  (1-\chi^*(p)p^{-s}). \tag{9.1}
\]
Meromorphic continuation preserves this finite-product identity.
For a prime in the product, $|\chi^*(p)|=1$, so a zero of the Euler
factor satisfies $p^{-\Re s}=1$, hence $\Re s=0$.
Therefore these factors create and cancel no open-strip zeros.
The principal character reduces to the primitive character of conductor
one, whose $L$-function is $\zeta(s)$; its pole at $s=1$ is not a zero.
This is exactly the induced-character convention in \sref{30006676}{1},
equations 25.15.1--25.15.5. Proving the assertion for primitive characters,
including conductor one, suffices for every character in the entry.

\paragraph{Attempt: an $L^2$ criterion for a holomorphic reciprocal.}
For a fixed primitive $\chi$ set
\[
 M_\chi(x)=\sum_{1\leq n\leq x}\mu(n)\chi(n),
\]
where $\mu(n)=0$ when a square prime factor divides $n$, and
$\mu(n)=(-1)^r$ when $n$ is a product of $r$ distinct primes.
Consider the following proposed estimate, separately for every $\chi$:
\[
 \text{for every }\varepsilon>0,\qquad
 I_{\chi,\varepsilon}:=
    \int_1^\infty |M_\chi(x)|^2x^{-2-\varepsilon}\,dx<\infty.
 \tag{9.2}
\]
Constants may depend on $\chi$ and $\varepsilon$. A uniform bound over
all conductors is \emph{not} added to this sufficient condition.

Assuming (9.2), define
\[
 H_\chi(s)=s\int_1^\infty M_\chi(x)x^{-s-1}\,dx
                  \quad(\Re s>1/2). \tag{9.3}
\]
To justify this definition, write $\sigma=\Re s$ and choose
$0<\varepsilon<2\sigma-1$. Cauchy--Schwarz gives
\[
 \int_1^\infty |M_\chi(x)|x^{-\sigma-1}\,dx
 \leq I_{\chi,\varepsilon}^{1/2}
       \left(\int_1^\infty x^{-2\sigma+\varepsilon}\,dx\right)^{1/2}
 =\frac{I_{\chi,\varepsilon}^{1/2}}{
                    \sqrt{2\sigma-\varepsilon-1}}. \tag{9.4}
\]
For a compact subset of the half-plane, choose a single $\varepsilon$
smaller than twice its distance from the boundary. The same argument,
with $(\log x)^{2j}$ in the second integral, bounds every derivative of
the integrand. Consequently (9.3) is holomorphic, by locally uniform
convergence and differentiation under the integral. There is no
unjustified interchange at $\Re s=1/2$; that boundary is not included.

In $\Re s>1$, the absolutely convergent Euler products give
\[
 H_\chi(s)=\sum_{n\geq1}\frac{\mu(n)\chi(n)}{n^s}
          =\frac1{L(s,\chi)}. \tag{9.5}
\]
The first equality is summation by parts; $M_\chi(x)=O(x)$ is enough
for its endpoint term to vanish there. Since $L$ has the continuation
specified in the original statement, meromorphic identity extends
$L(s,\chi)H_\chi(s)=1$ to $\Re s>1/2$. A zero of $L$ at a point where
it is holomorphic contradicts this identity. At $s=1$ for conductor one,
$H_\chi$ instead has a zero cancelling the pole, which is harmless.
The primitive functional equation relates a zero $\rho$ to a zero
$1-\rho$ of the conjugate character's completed function. Gamma factors
are finite and nonzero in the open strip. Applying the right-half-plane
argument to all primitive characters, including their conjugates,
therefore excludes left-half-strip zeros as well. \sref{30006676}{1}.
Together with (9.1), this proves the original GRH assertion under (9.2).

\paragraph{Attempt: a dyadic estimate would close the analytic step.}
It would suffice to show, for every $\varepsilon>0$ and $X\geq1$,
\[
 \int_X^{2X}|M_\chi(x)|^2\,dx
       \leq C_{\chi,\varepsilon}X^{2+\varepsilon/2}. \tag{9.6}
\]
Indeed the contribution of $[2^j,2^{j+1}]$ to (9.2) is at most
$C_{\chi,\varepsilon}2^{-j\varepsilon/2}$, a convergent geometric
series. This mean-square condition allows spikes and is not silently
replaced by a pointwise bound. The trivial estimate $|M_\chi(x)|\leq x$
gives a block integral of order $X^3$, which misses the needed exponent
by nearly one. The integer-squarefree weight $\mu$ is thus the substantive
cancellation issue; the Mellin analysis by itself does not improve it.

\paragraph{Stress test: cancellation of $\chi$ is not cancellation of $\mu\chi$.}
For a nonprincipal character modulo $q$, summing over a complete period
gives zero. Thus $\sum_{n\leq x}\chi(n)=O(q)$, yielding a direct
continuation formula for $L(s,\chi)$ on $\Re s>0$. There is no corresponding
periodicity for $\mu(n)\chi(n)$: divisibility by squares and parity of
the number of prime factors are not periodic modulo $q$. Replacing
$M_\chi$ by the ordinary character sum would prove continuation of
$L$, not continuation of its reciprocal.

A small analytic countermodel pinpoints that logical difference. The
Dirichlet polynomial
\[
 D(s)=1-\tfrac32\,2^{-s}
\]
has bounded coefficient partial sums and is entire, but vanishes at
$\beta=\log(3/2)/\log2$, with $1/2<\beta<1$. Its reciprocal has a pole
there. This polynomial is not asserted to be a Dirichlet $L$-function;
it disproves only the inference from bounded coefficient sums and
continuation to zero-freeness. Likewise, agreement of numerical partial
sums with square-root growth up to a finite cutoff would not establish
(9.2), whose integral runs to infinity.

\paragraph{Outcome.}
\textbf{Outcome: Conditional reduction.}
The attempt gives a fully justified sufficient mean-square criterion,
with explicit compact convergence, principal-character treatment, and
finite-Euler-factor bookkeeping. A proof of (9.2), or the stronger
dyadic bound (9.6), for every primitive character would settle the exact
entry. Neither estimate is established here, and the argument does not
claim a converse or historical novelty for the reciprocal-Mellin method.
The obstruction is arithmetic cancellation, not an omitted analytic
continuation step.
%% END RESEARCH INSERTION
\subsection{Sources}
\begin{itemize}\raggedright
\item[\textnormal{[S1]}]\hypertarget{source:30006676:1}{} NIST DLMF25.15, Dirichlet L-functions, equations25.15.1–4 and zeros; GRH convention in Kandhil–Languasco–Moree, Mathematische Annalen394:43(2026).
\url{https://dlmf.nist.gov/25.15}.
{\small\textit{Catalogue formulation source. Consulted 24 September 2026: Dirichlet characters, L-functions and zeros.}}
\item[\textnormal{[S2]}]\hypertarget{source:30006676:2}{} Large values of quadratic character sums — Dong, Wang, Wang, Zhang and Zhao.
\url{https://arxiv.org/abs/2608.15773}.
{\small\textit{Catalogue research/status reference; not a proof endorsement.}}
\item[\textnormal{[S3]}]\hypertarget{source:30006676:3}{} On the Extended, Generalized, and Grand Riemann Hypotheses: A Unified Approach Based on the General Properties of L-Functions — Weicun Zhang, v19.
\url{https://www.preprints.org/manuscript/202506.0481}.
{\small\textit{Catalogue research/status reference; not a proof endorsement.}}
\end{itemize}

%% END ORIGINAL SECTION CONTAINER
\end{document}
